\chapter{Files, Formats, and Data Safety}

MobileCAD has no cloud of its own: files are files, read and written through Android's system file picker (SAF),
and you decide where they live. This chapter explains the formats, the save workflow, automatic recovery, and
the usual pitfalls.

\section{The native document: FCStd}

\begin{itemize}
  \item \file{.FCStd} is FreeCAD's native document: it saves \textbf{the entire modelling history} (sketches, features,
        expressions, assemblies, drawings), stays small, and can be fully restored;
  \item Saving keeps a \file{.fcstd1} backup;
  \item The headless engine does not write \file{GuiDocument.xml} (interface and display information) by default;
        MobileCAD \textbf{preserves and writes back} colours, transparency, face colours, and visibility---so when you save
        the file back and open it on the desktop, the appearance is not lost;
  \item When you open a \file{.FCStd} that contains an assembly, the app switches to the ``Assembly'' page automatically.
\end{itemize}

\section{Import and export}

\subsection{Import}

\ui{File → Import} or \ui{Part → Import}; the format is detected by content sniffing:

\begin{center}
\small
\begin{tabularx}{\linewidth}{@{}lX@{}}
\toprule
\textbf{Format} & \textbf{Handled by} \\
\midrule
STEP / IGES & Part import (Import)\\
STL / PLY / OFF / AMF / 3MF & Mesh\\
BREP & Part\\
IFC & NativeIFC (BIM)\\
DXF / SVG / OBJ / DAE / CSG / 3DS / SHP, etc. & Importers registered by FreeCAD \\
FEM results (\file{.frd}, Z88, Fenics) & FEM result import \\
\bottomrule
\end{tabularx}
\end{center}

\subsection{Export}

\ui{File → Export}: \textbf{the extension is the format}---

\begin{itemize}
  \item Solids: STEP (\file{.step/.stp}), IGES, BREP, STL, OBJ, PLY, OFF, AMF, 3MF;
  \item Buildings: IFC, IFCJSON;
  \item 2D and drawings: DXF, SVG (flattened), PDF;
  \item Other: WebGL (\file{.html}), CSG/SCAD, JSON;
  \item Images: screenshots as PNG; drawing pages as PDF/SVG/DXF; spreadsheets as CSV/TSV; CAM G-code; robot \file{.src}.
\end{itemize}

\begin{warn}[Two things that do not work]
\begin{itemize}
  \item \textbf{DWG cannot be opened}: FreeCAD's DWG import relies on an external ODA/LibreDWG/QCAD converter,
        which the APK does not include; convert the file to DXF first.
  \item \textbf{glTF import may be unavailable}: the runtime lacks the corresponding import module; convert to
        STEP/STL first if you need it.
\end{itemize}
\end{warn}

\section{The system file picker (SAF)}

Every file operation goes through Android's SAF:

\begin{enumerate}
  \item When you trigger open, save, import, or export, the app first records the request (\ui{FileRequest}) and then
        brings up the system picker;
  \item \textbf{Open}: copies the selected file into the app cache and hands it to the engine; the app requests
        persistent read permission so that ``Recent files'' still work after a restart;
  \item \textbf{Save}: the engine writes a cache file first, then writes it back through an output stream to the URI you chose;
  \item The file picker does not filter extensions (it shows \file{*/*} only), so \textbf{always type the extension
        yourself when saving};
  \item If the file has no extension when you open it (for example a content provider returns \file{document:1234}),
        the app sniffs the content: STEP header, IGES, BREP, STL length check, \file{Document.xml} in FCStd, 3MF, SVG, DXF, OBJ.
\end{enumerate}

\begin{tip}[Two ways to save]
\begin{itemize}
  \item \menupath{File → Save}: \textbf{overwrites silently} when the location is known, without opening the picker;
  \item \menupath{File → Save as}: always opens the picker so you can specify a new location or name.
\end{itemize}
\end{tip}

\section{Recent files and other document commands}

\begin{itemize}
  \item \ui{Recent files}: keeps at most \textbf{8} entries in the app configuration; reopen documents quickly from here;
  \item \ui{Save a copy} (SaveCopy): exports a copy of the current document (keeping the view data as well);
  \item \ui{Merge projects} (MergeProjects): merges the objects of another \file{.FCStd} into the current document;
  \item \ui{Document information} (ProjectInfo): view or edit the title, author, company, licence, comments, and so on
        (6 of them are editable);
  \item \ui{Revert}: discards unsaved changes and reloads the version on disk.
\end{itemize}

\section{Autosave and crash recovery}

\begin{itemize}
  \item Every \textbf{60 seconds}, and whenever the app goes to the background, a modified document gets a recovery
        copy in the app's private directory \file{files/recovery/} (\file{recovery.FCStd} plus a JSON file recording
        the URI, name, and time);
  \item The recovery copy is deleted after a normal save;
  \item After the app crashes, the next start opens the recovery copy automatically and shows
        \ui{Recovered unsaved work from HH:mm}; save it to a real file immediately.
\end{itemize}

\section{Single document and memory}

MobileCAD is a \textbf{single-document interface}: opening a new file replaces the current view (the old document
stays in memory and may be reclaimed by the system). For large projects:
\begin{itemize}
  \item one \file{.FCStd} per project, split into several files if necessary and brought together with ``Merge projects'';
  \item save often, and use ``Save a copy'' at important milestones;
  \item save before switching apps.
\end{itemize}

\section{Units}

All lengths are stored internally in millimetres and masses in kilograms. \menupath{Preferences → Unit system}
changes only the display and input parsing (to inches, for example); it does not change the model's actual dimensions.
